\section{Glossary}
\label{app:glossary}

This glossary keeps the manuscript-specific vocabulary small. Standard
mechanistic interpretability terms are used in their usual sense; we define them
only where SRP uses a specialized readout version.

\subsection{Core Objects}

\begin{description}[leftmargin=0.28\linewidth,style=nextline]
    \item[Readout]
    The model's final unembedding or LM head map from hidden states to vocabulary
    logits. In this paper, readout refers to this final interface, not to the
    whole model.

    \item[Unembedding matrix, \(\WU\)]
    The matrix \(\WU\in\R^{|\mathcal{V}|\times d}\) whose rows map a hidden
    state to token logits. Row \(t\) is \(w_t^\top\), so token \(t\)'s logit is
    \(h^\top w_t\).

    \item[Selected readout direction and selected readout score]
    A fixed linear combination of LM-head rows,
    \(q_\alpha=\sum_t\alpha_t w_t\), and the scalar readout value it defines at
    hidden state \(h\), \(s_\alpha(h)=h^\top q_\alpha\).
\end{description}

\subsection{Common Readout Scores and Contrasts}

\begin{description}[leftmargin=0.28\linewidth,style=nextline]
    \item[Token logit]
    The special case \(q=w_A\), asking how high token \(A\)'s vocabulary logit is
    at hidden state \(h\).

    \item[Centered logit]
    A token logit after subtracting a shared reference row, such as
    \(h^\top(w_A-\bar{w}_{\mathcal{V}})\) or \(h^\top(w_A-\mu)\).

    \item[Pairwise logit difference]
    The contrast \(h^\top(w_A-w_B)\). Positive values favor \(A\); negative
    values favor \(B\). This is the main signed readout score in our
    experiments.

    \item[Vocabulary-mean contrast]
    A contrast against the vocabulary-mean row,
    \(w_A-\bar{w}_{\mathcal{V}}\), where
    \(\bar{w}_{\mathcal{V}}=|\mathcal{V}|^{-1}\sum_t w_t\).

    \item[Group contrast]
    A contrast between average rows for two token groups, such as an answer
    group versus an abstention group. This reduces dependence on one tokenizer
    row.

    \item[Top competitor margin]
    A contrast between the current top token and one or more nearby competitors,
    \(w_u-\sum_{r\in\mathcal{R}}\pi_r w_r\).
\end{description}

\subsection{Sparse Readout Prism}

\begin{description}[leftmargin=0.28\linewidth,style=nextline]
    \item[Sparse Readout Prism]
    The sparse feature decomposition introduced here: an SAE trained on
    unembedding rows factorizes the final readout into sparse code entries for
    token rows and shared readout directions, yielding a sparse feature basis
    for the readout and a reconstructed LM head.

    \item[Score decomposition]
    The downstream use of the SRP basis with a hidden state and selected readout
    direction to decompose a selected readout score into signed feature terms plus
    reconstruction error.

    \item[Readout SAE]
    The sparse autoencoder trained on unembedding rows. It reconstructs a row as
    \(\widehat{w}_t=\mu+\sum_i z_{t,i}d_i\).

    \item[SAE feature / decoder direction, \(d_i\)]
    A reusable direction learned by the readout SAE. Its projection at hidden
    state \(h\) is \(h^\top d_i\).

    \item[Sparse coefficient / code entry, \(z_{t,i}\)]
    The sparse code entry with which token row \(t\) uses SAE feature \(i\).

    \item[Shared offset, \(\mu\)]
    The common readout offset used by the SAE. For zero sum contrasts such
    as logit differences and group contrasts, its contribution cancels.

    \item[Native \(k\)]
    The active feature budget of the trained TopK SAE. Display cutoffs may hide
    small bars, but they do not change this reconstruction budget.

    \item[Projection profile]
    The ranked pattern of readout feature projections \(h^\top d_i\) before
    selecting a readout score.
\end{description}

\subsection{Decomposition and Measures}

\begin{description}[leftmargin=0.28\linewidth,style=nextline]
    \item[Direction coefficient, \(\beta_i(\alpha)\)]
    The amount the selected readout direction uses SAE feature \(i\):
    \(\beta_i(\alpha)=\sum_t\alpha_t z_{t,i}\).

    \item[Signed feature contribution, \(c_i(h,\alpha)\)]
    The signed term \(c_i(h,\alpha)=\beta_i(\alpha)h^\top d_i\). Positive terms
    raise the selected readout score; negative terms lower it.

    \item[Score decomposition at a hidden state]
    The decomposition obtained by combining a selected readout direction with the SRP
    readout basis at hidden state \(h\). It reports signed feature terms plus
    reconstruction error for the selected readout score.

    \item[Feature sum, \(s_{\mathrm{feat}}\)]
    The sum of all signed feature contributions, \(\sum_i c_i(h,\alpha)\), before
    adding offset or reconstruction error terms.

    \item[Reconstruction error, \(\epsilon\)]
    The difference between the exact selected readout score and the reconstructed score:
    \(s_{\mathrm{exact}}-s_{\mathrm{recon}}\).

    \item[Relative reconstruction error, \(\rho\)]
    The normalized error
    \(|s_{\mathrm{exact}}-s_{\mathrm{recon}}|/(|s_{\mathrm{exact}}|+\delta)\).
    Aggregate summaries use \(\delta=0.5\) logit.

    \item[Sign agreement / sign flip]
    Whether a reconstructed signed contrast has the same sign as the exact contrast.
    Sign flips are reported separately because they reverse which side of a
    contrast the sparse reconstruction supports.

    \item[Feature label]
    A short label derived from vocabulary rows most associated with an SAE
    feature. The feature id remains the stable object of analysis; the label is
    a readability aid.
\end{description}

\subsection{Evaluation Terms}

\begin{description}[leftmargin=0.28\linewidth,style=nextline]
    \item[Evaluation set]
    A fixed evaluation input: prompts, tokens or token groups, linear
    coefficients, and tokenization filters fixed before model evaluation. The
    readout SAE is not trained on these prompts.

    \item[Evaluation item]
    One evaluated score or contrast after tokenization filters. Prompt variants and
    score variants can share a base case, so evaluation items are tracked
    separately from prompts.

    \item[Base case]
    A cluster of variants derived from the same prompt/score template. Base
    cases are the bootstrap units for confidence intervals.

    \item[rowEV]
    Row-centered explained variance for reconstructed held-out unembedding
    rows. It is a capacity metric paired with replacement and score-reconstruction
    metrics for model decision evidence.

    \item[Top-1 agreement / top-5 overlap]
    Agreement between original and reconstructed LM head rankings on held-out
    hidden states. Top-1 compares the argmax token; top-5 overlap is the average
    overlap between the two top-five sets and is not top-five accuracy.

    \item[C4 replacement]
    The measurement that substitutes \(\widehat{\WU}\) for \(\WU\) on C4 hidden
    states and compares the resulting readout distribution with the original
    one using top-1 agreement, KL, and ordering agreement.

    \item[Softcap correct]
    A Gemma evaluation that applies the model's monotone final-logit softcap
    before computing top-1 agreement or KL, matching the model's actual output
    path.

    \item[Strict budget / high fidelity settings]
    Two native-\(k\) operating regimes for readout SAEs. The strict budget
    setting uses native \(k=128\) for fixed budget comparability; the
    high fidelity setting uses native \(k=256\) when readout fidelity is the
    primary constraint.

    \item[Sparse code usage]
    Statistics describing how often learned SAE features are active in
    evaluation. A rare feature rate is a usage statistic, not a readout fidelity
    metric.

    \item[Dead / rare feature]
    A usage statistic for learned SAE features. Dead features do not activate
    in evaluation; rare features activate below the reporting threshold used in
    Appendix~\ref{app:readout-factorizer-sweeps}.
\end{description}

\subsection{Attribution and Controls}

\begin{description}[leftmargin=0.28\linewidth,style=nextline]
    \item[Direct logit attribution, DLA]
    Component attribution that projects residual stream components onto the
    readout direction for a selected readout score. We use DLA as additive attribution
    unless paired with interventions.

    \item[Feature-resolved DLA]
    The combined view that splits DLA component contributions across the
    readout features for a realized forward pass.

    \item[Additivity error]
    The gap between the exact selected readout score and the sum of component attribution
    terms. It is reported separately from the readout SAE reconstruction error.

    \item[Token audit item]
    A tokenizer or row factor, such as token length, frequency, row norm,
    capitalization, whitespace, punctuation, special token status, or
    tokenization collision.

    \item[Tokenization collision]
    A case where tokenization makes the intended lexical contrast brittle or
    ambiguous. Such cases are handled through token audits or converted
    to explicit group contrasts.
\end{description}
